\documentclass[10pt,a4paper]{amsart}
\usepackage{amsmath,amssymb,amsthm,graphicx}
\usepackage{hyperref}
% Minimal class-compatibility preamble for a reader excerpt of
%   Bhimarthi Ravinder, "Demazure Modules, Chari-Venkatesh Modules and Fusion Products",
%   SIGMA 10 (2014), 110, DOI 10.3842/SIGMA.2014.110.
%
% The author's pinned TeX file declares no theorem environments of its own: they are
% supplied by the SIGMA publisher class sigma.cls, which is not installed here.  The
% block below is that class's theorem-environment block (same environment names and
% styles), copied mechanically so that the retained author text compiles unchanged.
% No mathematical macro or statement is introduced by this file.
\newtheorem{theorem}{Theorem}
\newtheorem*{theorem*}{Theorem}
\newtheorem{lemma}{Lemma}
\newtheorem*{lemma*}{Lemma}
\newtheorem{corollary}{Corollary}
\newtheorem*{corollary*}{Corollary}
\newtheorem{proposition}{Proposition}
\newtheorem*{proposition*}{Proposition}
\newtheorem{conjecture}{Conjecture}
\newtheorem*{conjecture*}{Conjecture}
{\theoremstyle{definition} \newtheorem{definition}{Definition}
\newtheorem*{definition*}{Definition}
\newtheorem{example}{Example}
\newtheorem*{example*}{Example}
\newtheorem{remark}{Remark}
\newtheorem*{remark*}{Remark}
\newtheorem{note}{Note}
\newtheorem*{note*}{Note}
}
\DeclareMathOperator{\ev}{ev}
\DeclareMathOperator{\gr}{gr}
\DeclareMathOperator{\spn}{span}
\title{Explicit defining relations for the level-one Demazure module of a finite-dimensional complex simple Lie algebra}
\author{Bhimarthi Ravinder}
\date{doi:10.3842/SIGMA.2014.110; CC BY 4.0}
\begin{document}
\maketitle
\noindent\textit{Editorial note.} All mathematical text is the author's own, copied from the fixed publisher source; no statement or proof step was rewritten.\par\medskip
\noindent\textit{Editorial note.} The author's file declares no theorem environments of its own: they come from the SIGMA publisher class, which is not installed for this excerpt. The minimal preamble therefore reproduces that class's theorem-environment block (same environment names and styles); the author's own operator macros are taken from source lines 9-11.\par\medskip
\noindent\textit{Editorial note.} The bibliography is the author's own thebibliography, copied verbatim from source lines 875-924, so that the one citation occurring in the context block resolves. Environment numbers printed here are local to the excerpt; the source statement is the proposition labelled WvsD at source line 396 (Section 3).\par\medskip
\noindent\textit{Allowed standard prerequisites (editorial).}\begin{itemize}
\item Root systems and structure theory of finite-dimensional complex simple Lie algebras. Roots, coroots, the highest root theta, the invariant form normalized on long roots and d\_alpha = 2/(alpha|alpha) are assumed. The dichotomy (theta|alpha) in \{0,1\} for alpha in R\textasciicircum{}+ different from theta, its consequence for <k theta, alpha\textasciicircum{}vee>, and the short/long root alternatives are part of the retained proof.
\item Local Weyl modules for the current algebra g[t] and their defining relations. The excerpt carries the author's definition of W(lambda), its grading, and the relation (w2') that follows from Garland's identity. Comparing those relations with the defining families of the level-one Demazure quotient is the retained proof.
\end{itemize}\medskip
\begin{lemma}\label{garland}
Given $s\in\mathbb{N}$, $r\in\mathbb{Z}_{\geq0}$ and $\alpha\in R^{+}$, we have
\begin{gather*}
(x^{+}_{\alpha} \otimes t)^{(s)}(x^{-}_{\alpha} \otimes
1)^{(s+r)}-(-1)^s\mathbf{x}^{-}_{\alpha}(r,s)\in\mathbf{U}(\mathfrak{g}[t])\mathfrak{n}^{+}[t]\bigoplus
\mathbf{U}(\mathfrak{n}^{-}[t])\mathfrak{h}[t]_{+}.
\end{gather*}
\end{lemma}

\begin{definition}
Given $\lambda \in P^+$, the local Weyl module $W(\lambda)$ is the cyclic $\mathfrak{g}[t]$-module generated by an
element $w_\lambda$, with following def\/ining relations:
\begin{gather}
\mathfrak{n}^{+}[t]  w_\lambda=0,
\qquad
(h \otimes t^s)  w_\lambda = \langle \lambda, h \rangle \delta_{s,0} w_\lambda=0,
\qquad
s\geq0,
\qquad
h\in\mathfrak{h},
%\label{w1}
\nonumber
\\
(x^{-}_\alpha \otimes 1)^{\langle \lambda,
\alpha^{\vee} \rangle + 1}w_\lambda=0,
\qquad
\alpha \in R^{+}.
\label{w2}
\end{gather}
\end{definition}
\noindent We note that the relation~\eqref{w2} implies
\begin{gather}
\label{w2'}
\big(x^{-}_\alpha \otimes t^{\langle \lambda,
\alpha^{\vee} \rangle}\big)  w_\lambda=0,
\qquad
\alpha \in R^{+},
\end{gather}
which is easy to see from Lemma~\ref{garland}.

\begin{definition}
\textup{(see~\cite[Corollary 3.5]{CV})} The level one Demazure module $D(1,\lambda)$ is the graded quotient of
$W(\lambda)$ by the submodule generated by the union of the following two sets:
\begin{gather}
\big\{(x^{-}_{\alpha} \otimes t^{s_{\alpha}})  w_{\lambda}: \alpha\in R^{+}~\textup{such that}~d_{\alpha} > 1\big\},
\label{dm1}
\\
\big\{\big(x^{-}_{\alpha} \otimes t^{s_{\alpha}-1}\big)^2  w_{\lambda}: \alpha\in R^{+}~\textup{such that}~d_\alpha =3~\textup{and}~m_\alpha=1\big\}.
\label{dm2}
\end{gather}
In particular, for $\mathfrak{g}$ simply laced, $D(1,\lambda)\cong_{\mathfrak{g}[t]} W(\lambda)$. We denote~by
$\overline{w}_\lambda$, the image of $w_\lambda$ in~$D(1,\lambda)$.
\end{definition}

\begin{proposition}
\label{WvsD}
Given $k\geq 1$, the level~$1$ Demazure module $D(1,k\theta)$ is the graded $\mathfrak{g}[t]$-module generated by an
element $\overline{w}_{k\theta}$, with the following defining relations:
\begin{gather*}
\mathfrak{n}^{+}[t]\, \overline{w}_{k\theta}=0,
\qquad
(h \otimes t^s) \overline{w}_{k\theta} = \langle k\theta, h\rangle \delta_{s,0} \overline{w}_{k\theta},
\qquad
s\geq0,
\quad
h\in\mathfrak{h},
\\
(x^{-}_{\alpha}\otimes 1) \overline{w}_{k\theta}=0,
\qquad
\alpha\in R^+,
\quad
(\theta|\alpha)=0,
\\
(x^{-}_{\alpha} \otimes 1)^{kd_{\alpha}+1} \overline{w}_{k\theta}=0,
\qquad
\big(x^{-}_{\alpha} \otimes t^k\big) \overline{w}_{k\theta}=0,
\qquad
\alpha\in R^+,
\quad
(\theta|\alpha)=1,
\\
(x^{-}_{\theta} \otimes 1)^{2k+1} \overline{w}_{k\theta}=0.
\end{gather*}
\end{proposition}

\begin{proof}
Observe that, from the abstract theory of root systems $(\theta|\alpha)= 0$ or~1,
$\forall\, \alpha \in R^{+}\setminus \{\theta\}$.
This implies that $\langle k\theta, \alpha^{\vee}\rangle = 0$ or $kd_{\alpha}$, $\forall\, \alpha \in R^{+}\setminus \{\theta\}$.
Hence the relations~\eqref{dm2} do not occur in $D(1,k\theta)$ and the relations~\eqref{dm1} are
\begin{gather*}
\big(x^{-}_{\alpha} \otimes t^k\big) \overline{w}_{k\theta}=0,
\qquad
\alpha\in R^+,
\quad
\alpha~~\text{short},
\quad
(\theta|\alpha)=1.
\end{gather*}
For a~long root $\alpha\in R^+$ with $(\theta|\alpha)=1$, by~\eqref{w2'} it follows that $(x^{-}_{\alpha} \otimes t^k) \overline{w}_{k\theta}=0$.
Now the other relations are precisely the def\/ining relations of $W(k\theta)$. This proves Proposition~\ref{WvsD}.
\end{proof}

% Bibliography of the delivered excerpt: the author's own thebibliography,
% copied verbatim from pinned-original-source.tex lines 875-924.
% The SIGMA-class-only end-of-article token (line 924) is omitted because the
% excerpt is typeset with a class-compatibility preamble instead of sigma.cls.
\begin{thebibliography}{99}
\footnotesize \itemsep=0pt

\bibitem{CFK}
Chari V., Fourier G., Khandai T., A categorical approach to {W}eyl modules,
  \href{http://dx.doi.org/10.1007/s00031-010-9090-9}{\textit{Transform. Groups}} \textbf{15} (2010), 517--549, \href{http://arxiv.org/abs/0906.2014}{arXiv:0906.2014}.

\bibitem{CP}
Chari V., Pressley A., Weyl modules for classical and quantum af\/fine algebras,
  \href{http://dx.doi.org/10.1090/S1088-4165-01-00115-7}{\textit{Represent. Theory}} \textbf{5} (2001), 191--223,
  \href{http://arxiv.org/abs/math.QA/0004174}{math.QA/0004174}.

\bibitem{CSVW}
Chari V., Shereen P., Venkatesh R., Wand J., A Steinberg type decomposition
  theorem for higher level Demazure modules, \href{http://arxiv.org/abs/1408.4090}{arXiv:1408.4090}.

\bibitem{CV}
Chari V., Venkatesh R., Demazure modules, fusion products, and $Q$-systems,
  \href{http://dx.doi.org/10.1007/s00220-014-2175-x}{\textit{Comm. Math. Phys.}}, {t}o appear, \href{http://arxiv.org/abs/1305.2523}{arXiv:1305.2523}.

\bibitem{FL}
Feigin B., Loktev S., On generalized {K}ostka polynomials and the quantum
  {V}erlinde rule, in Dif\/ferential Topology, Inf\/inite-Dimensional {L}ie
  Algebras, and Applications, \textit{Amer. Math. Soc. Transl. Ser.~2}, Vol.~194, Amer. Math. Soc., Providence, RI, 1999, 61--79,
  \href{http://arxiv.org/abs/math.QA/9812093}{math.QA/9812093}.

\bibitem{F}
Feigin E., The {PBW} f\/iltration, {D}emazure modules and toroidal current
  algebras, \href{http://dx.doi.org/10.3842/SIGMA.2008.070}{\textit{SIGMA}} \textbf{4} (2008), 070, 21~pages,
  \href{http://arxiv.org/abs/0806.4851}{arXiv:0806.4851}.

\bibitem{FoL}
Fourier G., Littelmann P., Weyl modules, {D}emazure modules, {KR}-modules,
  crystals, fusion products and limit constructions, \href{http://dx.doi.org/10.1016/j.aim.2006.09.002}{\textit{Adv. Math.}}
  \textbf{211} (2007), 566--593, \href{http://arxiv.org/abs/math.RT/0509276}{math.RT/0509276}.

\bibitem{G}
Garland H., The arithmetic theory of loop algebras, \href{http://dx.doi.org/10.1016/0021-8693(78)90294-6}{\textit{J.~Algebra}}
  \textbf{53} (1978), 480--551.

\bibitem{Naoi}
Naoi K., Weyl modules, {D}emazure modules and f\/inite crystals for non-simply
  laced type, \href{http://dx.doi.org/10.1016/j.aim.2011.10.005}{\textit{Adv. Math.}} \textbf{229} (2012), 875--934,
  \href{http://arxiv.org/abs/1012.5480}{arXiv:1012.5480}.

\bibitem{V}
Venkatesh R., Fusion product structure of Demazure modules, \href{http://dx.doi.org/10.1007/s10468-014-9495-6}{\textit{Algebr.
  Represent. Theory}}, {t}o appear, \href{http://arxiv.org/abs/1311.2224}{arXiv:1311.2224}.

\end{thebibliography}
\end{document}
